% --- Plik: chapter5_conclusions.tex ---

\chapter{Podsumowanie}
\label{ch:podsumowanie} % To jest unikalny znacznik tego rozdziału

Niniejszy rozdział stanowi formalne zamknięcie merytorycznej części pracy dyplomowej. Jego celem jest zwięzła synteza uzyskanych rezultatów oraz globalna ocena stopnia realizacji zadania inżynierskiego.

W rozdziale tym należy dokonać całościowej oceny zrealizowanych prac inżynierskich, opierając strukturę wniosków na następujących kluczowych aspektach:
\begin{itemize}
	\item jednoznacznej weryfikacji stopnia realizacji celów oraz tez sformułowanych we wstępie pracy,
	
	\item zwięzłym omówieniu najważniejszych osiągnięć o charakterze konstrukcyjnym, programistycznym bądź eksperymentalnym,
	
	\item obiektywnej ocenie ograniczeń wypracowanego rozwiązania oraz wskazaniu potencjalnych kierunków jego dalszej optymalizacji lub rozbudowy.
\end{itemize}
\noindent \textbf{Stylistyka:} wnioski powinny być sformułowane w sposób konkretny i zwięzły. Należy unikać dosłownego powtarzania zdań z poprzednich rozdziałów, skupiając się na ogólnej konkluzji z całej pracy.

\vspace{0.5cm}

\noindent \textbf{Więcej informacji na temat systemu \LaTeX{}:}
\begin{itemize}
	\item \url{https://www.overleaf.com/learn} -- bardzo przystępny, oficjalny tutorial i baza wiedzy na stronie platformy Overleaf,
	
	\item \url{https://www.latex-project.org/} -- oficjalna strona domowa projektu \LaTeX{}, zawierająca aktualności i dokumentację standardów,
	
	\item \url{https://www.tug.org/begin.html} -- prowadzony przez \textit{\TeX{} Users Group} obszerny zbiór odnośników do materiałów dla osób początkujących,
	
	\item \url{https://www.texstudio.org/} -- oficjalna strona domowa edytora TeXstudio, z której należy pobierać najnowsze, stabilne wersje programu,
	
	\item \url{https://www.xm1math.net/texmaker/} -- oficjalna strona domowa edytora Texmaker, z której można pobrać program oraz uzyskać dostęp do jego podręcznika użytkownika,

	\item \url{https://texstudio-org.github.io/} -- oficjalny, szczegółowy podręcznik użytkownika (\textit{User Manual}) wyjaśniający zaawansowane funkcje edytora TeXstudio,
	
	\item \url{https://tex.stackexchange.com/} -- największe na świecie forum społecznościowe typu Q\&A (pytania i odpowiedzi), na którym można znaleźć rozwiązanie niemal każdego problemu technicznego i błędu kompilacji.
\end{itemize}
